\section*{Chapter \ref{ch:m1:the-buy-side} --- The Buy Side}

\begin{solution}{exo:m1:the-buy-side:1}
Active returns $+1.1$, $-0.6$, $+0.8$, $+1.5$; mean $0.70\%$; sample standard
deviation $0.91\%$; $\mathrm{IR} = 0.77$ --- on four observations, which
\cref{prop:m1:the-buy-side:years} says is worth nothing yet.
\end{solution}

\begin{solution}{exo:m1:the-buy-side:2}
$4/0.4^2 = 25$ years; $4/1.5^2 = 1.8$ years; $4/36$ of a year, 28 trading
days.
\end{solution}

\begin{solution}{exo:m1:the-buy-side:3}
Index fund: $50 \times 10^9 \times 0.0004 = \$20$ million. \omterm{def:m1:what-a-trading-firm-does:hedge-fund}{Hedge fund}:
management \$40 million, performance $0.2 \times (200 - 40) = \$32$ million,
\$72 million in total on a twenty-fifth of the assets.
\end{solution}

\begin{solution}{exo:m1:the-buy-side:4}
Positions above 5\% total $9.5 + 18 + 11 = 38.5\% \le 40\%$ and none exceeds
10\%: compliant. After the rise the stock weighs $12.35/102.85 = 12.0\%$:
above 10\%. The manager must sell about a sixth of the position, whatever it
thinks of the company --- an order with no information in it.
\end{solution}

\begin{solution}{exo:m1:the-buy-side:5}
$\mathrm{TE} = \sqrt{100 \times (0.005\times0.30)^2} = 1.5\%$. For 3\%, double
the active weights: 1.0\% per position.
\end{solution}

\begin{solution}{exo:m1:the-buy-side:6}
Year 1: management 2.00, $N' = 128.00$, performance $0.2\times28 = 5.60$,
close 122.40 (per 100). Year 2: management 2.45, $N' = 95.47$, below the mark
of 122.40: no \omterm{def:m1:the-buy-side:fees}{performance fee}, close 95.47. Year 3: management 1.91, $N' =
103.11$, still below 122.40: no fee, close 103.11.
\end{solution}

\begin{solution}{exo:m1:the-buy-side:7}
``2 and 20'': fees 35.9 per 100 invested, final value 145.2. ``1 and 30'':
fees 33.2, final value 149.8. The investor who believes in the manager
prefers ``1 and 30'': it pays less in total here and pays mostly when it has
been made richer. The manager may prefer ``2 and 20'' because two thirds of
its income (23.4 of 35.9) then does not depend on performance at all, and
pays the salaries in the seven years spent under the mark.
\end{solution}

\begin{solution}{exo:m1:the-buy-side:8}
With no skill, a five-year mean active return is $N(0, 4\%/\sqrt5) = N(0,
1.79\%)$. The twentieth of 400 is the 95th percentile, $1.645 \times 1.79 =
2.9\%$ a year, an apparent \omterm{def:m1:the-buy-side:ir}{information ratio} of $2.9/4 = 0.74$; the top
handful exceed 1. The flaw is selection: the ranking guarantees impressive
records among pure noise, and these records have no bearing on the next five
years. The honest test is the performance of the selected funds
\emph{after} selection.
\end{solution}

\begin{solution}{pb:m1:the-buy-side:1}
Per \$100 invested.
\textbf{1.} Management 2.00; $N' = 118.00$; performance $0.2 \times 18 = 3.60$;
close 114.40; mark 114.40.
\textbf{2.} Management 2.29; $N' = 114.40\times0.85 - 2.29 = 94.95$; below
the mark, so no \omterm{def:m1:the-buy-side:fees}{performance fee}; mark stays 114.40.
\textbf{3.} Management 1.90; $N' = 102.55 < 114.40$: no fee.
\textbf{4.} The investor has $+2.55$; the manager has been paid $9.79$.
\textbf{5.} Years 1, 4 and 7.
\textbf{6.} Management 23.4; performance 12.4; total \$35.9 million.
\textbf{7.} \$145.2 million; 3.8\% a year.
\textbf{8.} \$192.8 million; 6.8\% a year.
\textbf{9.} $92.8 - 35.9 - 45.2 = 11.7$: the growth that the fees would have
earned had they stayed invested. Fees are paid early and compound against
the investor.
\textbf{10.} The manager received 39\% of the gross profit; the investor
kept 49\%; 12\% was the lost compounding.
\textbf{11.} $100 \times 1.0595^{10} = \$178.2$ million against \$145.2
million.
\textbf{12.} Diversification: a stream uncorrelated with its equities lowers
the volatility of the whole endowment, and it lost less in years 2 and 8 if
those were bad equity years. Whether that is worth 33 percentage points over
a decade is the board's question.
\textbf{13.} Mean gross return 8\%; Sharpe ratio $0.08/0.17 = 0.47$;
$4/0.47^2 = 18$ years.
\textbf{14.} Net value 116.8 against a mark of 149.7: a rise of 28\%.
\textbf{15.} (a) The option is far out of the money and gains value with
volatility; extra risk costs the manager nothing more if it fails, while
the investors bear the loss. (b) A new fund starts with a mark at par: the
manager escapes the 28\% gap, and the old investors lose the ``free'' years
they had paid for with their losses.
\textbf{16.} Its mark is its subscription value. In year 9: management 2.00,
$N' = 113.00$, performance $0.2 \times 13 = 2.60$ per 100 --- while the
first investor pays no \omterm{def:m1:the-buy-side:fees}{performance fee} that year.
\textbf{17.} The two investors hold the same portfolio and owe different
fees; a single fund-level mark would either charge the first investor
twice for the same gains or let the second ride free.
\textbf{18.} About 2.4\% a year at 20\% volatility against 1.8\% at zero
volatility: some 0.6\% a year for bearing risk with no skill.
\textbf{19.} \textbf{39\%.}
\textbf{20.} Neither number first: the terms that shape behaviour --- the
\omterm{def:m1:the-buy-side:fees}{high-water mark} and its reset conditions, a hurdle, the manager's own
capital in the fund and risk limits --- and then the \omterm{def:m1:the-buy-side:fees}{management fee}, which
is the part paid for nothing.
\end{solution}

\begin{solution}{iq:m1:the-buy-side:1}
\omterm{def:m1:the-buy-side:benchmark}{Tracking error} is the volatility of the difference between the portfolio's
return and the \omterm{def:m1:the-buy-side:benchmark}{benchmark}'s. A 15\% position in a stock that is 2\% of the
index is a 13\% active weight: with 30\% specific volatility it alone adds
about 4\% of \omterm{def:m1:the-buy-side:benchmark}{tracking error}, probably more than the \omterm{def:m1:the-buy-side:mandate}{mandate} allows, and one
bad quarter in that stock loses the \omterm{def:m1:the-buy-side:mandate}{mandate} however right the view is in
the long run. For a public fund a 10\% legal cap applies anyway.
\iqlookfor{active weight, not absolute weight, and career risk as a real
constraint.}
\end{solution}

\begin{solution}{iq:m1:the-buy-side:2}
$t = \mathrm{SR}\sqrt{T}$, so $t = 2$ needs 4 years, under ideal
conditions. In practice more: returns are not independent or normal, and if
the strategy was chosen among many the threshold must be far higher than 2.
\iqlookfor{the formula immediately, then the multiple-testing caveat
unprompted.}
\end{solution}

\begin{solution}{iq:m1:the-buy-side:3}
No. The index fund trades because of flows or an index change: no
information about the stock, so the mechanical discount applies
(\cref{prop:m1:the-sell-side:discount}). The \omterm{def:m1:what-a-trading-firm-does:hedge-fund}{hedge fund} may be selling on a
view or ahead of news: add an adverse-selection charge estimated from what
happened after its previous orders. Relationship value can cut the other
way.
\iqlookfor{pricing the client, not just the order.}
\end{solution}

\begin{solution}{iq:m1:the-buy-side:4}
``You pay the manager a share of profits, but only on new profits: if the
fund falls, it must climb back to its old peak before you pay that share
again.'' For the manager it is a call option struck at the peak: near the
mark it behaves normally; far below, the option is nearly worthless and
gains from volatility, so the manager is tempted to gamble or to close the
fund and reopen with a fresh mark.
\iqlookfor{a plain first sentence, then the option argument.}
\end{solution}

\begin{solution}{iq:m1:the-buy-side:5}
Index reconstitutions: every tracker must buy additions and sell deletions
at the effective close; whoever buys after the announcement and sells into
that close earns the concession, at the trackers' expense. Periodic
rebalancing and leveraged-fund resets: month-end and end-of-day flows of
known sign and size; liquidity providers position ahead of them. Also
closing auctions themselves, which concentrate this volume.
\iqlookfor{flows that are predictable in sign, size \emph{and} timing.}
\end{solution}

\begin{solution}{iq:m1:the-buy-side:6}
After three years an unskilled strategy's estimated Sharpe ratio is about
$N(0, 1/\sqrt3) = N(0, 0.58)$. The maximum of 1\,000 independent normals is
about 3.2 standard deviations: a Sharpe ratio near 1.9. A backtested Sharpe
of 2 chosen among a thousand trials is what nothing looks like.
\iqlookfor{the expected maximum of $n$ normals, roughly $\sqrt{2\ln n}$,
and the conclusion about research practice.}
\end{solution}
